\documentclass{article}
\usepackage[T1]{fontenc}
\usepackage{lmodern}
\usepackage{graphicx}
\usepackage{amsmath,amssymb}
\usepackage[a4paper,margin=2cm]{geometry}
\usepackage[absolute,overlay]{textpos}
\setlength{\TPHorizModule}{1mm}
\setlength{\TPVertModule}{1mm}
\usepackage[indonesian]{babel}
\usepackage{hyperref}
\setlength{\emergencystretch}{2em}
% unit: Halaman 31

\begin{document}

% === Halaman 31 (llm) ===

\section*{S-box}

Elemen-elemen di dalam $S\text{-}box$ terlihat seperti diisi secara acak, namun sebenarnya ia dihasilkan dari proses perhitungan sebagai berikut:

\begin{enumerate}
    \item Inisialisasi $S\text{-}box$ dengan nilai yang menaik dari baris ke baris. Baris ke-0 diisi dengan nilai 00, 01, 02, ..., 0F. Baris ke-1 diisi dengan nilai 10, 11, 12, ..., 1F, dan seterusnya untuk baris-baris lain.
    \item Untuk setiap nilai pada baris $y$ kolom $x$, tentukan balikannya dalam $\text{GF}(2^8)$. Nilai 00 dipetakan ke dirinya sendiri. Penjelasan tentang $\text{GF}(2^8)$ lihat pada lampiran.
    \item Hasil dari langkah 2 dikonversi ke vektor kolom bit $(b_0, b_1, ..., b_8)^\top$.
\end{enumerate}

\end{document}
